\chapter{Related Work}
This chapter follows the argument of the introduction. We begin with how
SLR has treated signer variation, as an evaluation protocol and as
something to suppress (\S\ref{sec:related-signer-independent}). We then
turn to appearance bias itself, which is well studied in action
recognition and barely touched in sign language, and to why the action
recognition remedies do not transfer
(\S\ref{sec:related-appearance-bias}). We then trace how benchmarks became
more diverse and what that did to modality preference
(\S\ref{sec:related-work-benchmarks}), examine the resulting move from RGB
to pose and what it costs (\S\ref{sec:related-work-modality}), and close
with synthetic data as the intervention family this thesis draws on
(\S\ref{sec:related-work-synthetic}).

\section{Signer-independent SLR}
\label{sec:related-signer-independent}
Signer variation is a long-standing generalisation axis in SLR, and the
field has responded on two fronts: evaluation protocols that withhold
signers, and methods that suppress signer information.

On the protocol side, signer-independent evaluation is well established.
The ChaLearn LAP challenge on AUTSL made large-scale signer-independent
isolated recognition a community benchmark (226 signs, 43 signers,
signer-disjoint splits); its winning entries relied heavily on pose
estimation, graph convolutions, and
ensembling~\cite{sincanChaLearnLAPSigner2021}. PHOENIX-2014 ships a
dedicated signer-independent split alongside its default protocol, used
by signer-invariant CSLR work~\cite{zuoImprovingContinuousSign2024},
and ASL Citizen is signer-disjoint by construction~\cite{aslcitizen}.
Recent challenges push the setting further: the Cross-View Isolated SLR
challenge adds camera viewpoint to the held-out axes~\cite{CV-ISLR},
and the SignEval 2025 challenge makes signer-independent CSLR on Isharah
a dedicated track, distributing only pose features extracted from the
RGB recordings out of concern for signer overfitting and
privacy~\cite{luqmanSignEval2025Challenge2025a,
alyamiIsharahLargeScaleMultiScene2026b}, a design choice we return to in
\S\ref{sec:related-work-modality}.

On the method side, three families recur. Adversarial approaches train
the encoder so that signer identity cannot be recovered from it:
Ferreira et al.\ apply signer-adversarial training to isolated
recognition~\cite{Ferreira2021}, and Zuo and Mak add a signer removal
module that disentangles signer statistics from the recognition features
of a CSLR model~\cite{zuoImprovingContinuousSign2024}. Disentanglement
approaches separate signer style from sign content explicitly, as in
contrastive disentangled meta-learning for signer-independent
translation~\cite{jinContrastiveDisentangledMetaLearning2021}. Data-side
approaches raise the signer diversity the model sees, with more real
signers or with synthetic signer variety such as SDDA's
diffusion-generated signers (\S\ref{sec:related-work-synthetic}).
Pose-centric modelling, finally, sidesteps the pixel appearance
altogether~\cite{kratimenosIndependentSignLanguage2020,
tranRegionAwarePoseModeling2025}, at the cost examined in
\S\ref{sec:related-work-modality}.

What these approaches share is that appearance is handled wholesale: held
out with the signer, removed adversarially, averaged over by adding
signers, or discarded along with the pixels. None isolates which visual
properties carry the signal, or tests whether changing exactly those,
with the signing held fixed, removes the shortcut.
\S\ref{sec:related-appearance-bias} takes up that question.
\section{Addressing appearance bias}
\label{sec:related-appearance-bias}
That action-recognition benchmarks reward static appearance cues is an
empirical finding rather than an assumption: RESOUND quantifies the
representation bias of action datasets and shows the standard ones are
skewed towards static objects, scenes and
people~\cite{liRESOUNDActionRecognition2018}. RESOUND both resamples
existing corpora to reduce that bias and collects Diving48, a
fine-grained benchmark designed to remove it. The effect is severe enough
that actions stay recognisable once the actor is gone: masking out
detected people in UCF101 still leaves a two-stream network at $53.1\%$,
against $84.3\%$ on the unmodified
videos~\cite{heHumanActionRecognition2016}.

Crucially for this thesis, the bias is a property of the training data
rather than of the architecture.
\citet{byvshevAre3DConvolutional2022a} measure how much temporal
structure a 3D convolutional network's filters retain and find it
collapsing with depth in Kinetics-trained models, then retrain the same
architectures on synthetic data in which motion and appearance are
decoupled by construction; there the collapse disappears, and they
conclude that the architecture is not inherently appearance-biased.
\citet{kowalQuantifyingLearningStatic2025} reach a compatible conclusion
with a representation-level probe, pairing restyled videos (dynamics
fixed, appearance changed) against frame-shuffled ones (appearance fixed,
dynamics destroyed) to score how static each channel is, and find most
models static-biased. If the bias is inherited from data, it can be
attacked by reshaping data.

Remedies fall into two families. The first restructures the training
signal so that positives share motion while appearance varies. Motion
Coherent Augmentation permutes a clip's RGB channels and blends the
result back into the original, a cheap hue shift that leaves saturation
and value untouched, and pairs it with a loss aligning the augmented
clip's prediction to the original clip's rather than to the
label~\cite{zhangDontJudgeLook2024}. Tubelet-contrastive self-supervision
goes further, building positive pairs that share only local motion by
pasting synthetic moving tubelets into different videos, so that
appearance cannot solve the pretext
task~\cite{thokerTubeletContrastiveSelfSupervisionVideoEfficient2023b}.
The second family suppresses the static signal directly. StillMix uses a
2D network to flag frames from which the action is already predictable,
then tiles such a frame into a motionless clip and blends it into
training video~\cite{liMitigatingEvaluatingStatic2023a}; ALBAR repeats a
single frame into a motionless clip, passes it through the same encoder
as the real clip, and penalises any confident prediction from it with a
negative cross-entropy and an entropy
term~\cite{fioresiALBARAdversarialLearning2025}; SMAD combines
multi-aspect adversarial training with spatial actionness
reweighting~\cite{duanMitigatingRepresentationBias2023}; and
\citet{rastegarBackgroundNoMore2024} treat the background as a confounder
and intervene causally at both training and test time.

Two of these lines already reach past the background to the actor:
StillMix demonstrates that the foreground is itself a source of static
bias, and ALBAR names clothing, skin colour and facial hair explicitly.
Their mechanism, however, is to \emph{suppress} the static channel, by
masking it, blending a motionless clip into it, or penalising confident
predictions made from it. That is not available in signing, where the
same channel carries linguistic content: facial expression, mouthing and
fine handshape are part of what the label depends on. The intervention
must therefore change appearance while preserving motion, rather than
removing appearance information, which is what makes a counterfactual
rather than a suppression the natural instrument here.

Within sign language, the nearest work targets demographic disparity and
research practice rather than the appearance content of the training
distribution. Atwell et al.\ study and mitigate demographic bias on ASL
Citizen, reducing performance disparities without a loss of
accuracy~\cite{atwellStudyingMitigatingBiases2024}, and a Deaf-led review
documents systemic biases in how sign language AI research selects
datasets, labels and modelling
assumptions~\cite{desaiSystemicBiasesSign2024}.
SignRep~\cite{wongSignRepEnhancingSelfSupervised2025} intervenes on the
representation instead, suppressing appearance with an adversarial style
loss during self-supervised pretraining
(\S\ref{sec:related-work-modality}). The closest data-side intervention
is SDDA~\cite{fuSignerDiversitydrivenData2024b}, treated with the other
synthetic-data work in \S\ref{sec:related-work-synthetic}.

Editing the appearance of sign video has precedent, with a different
goal: diffusion-based editing has been used to \emph{anonymise} sign
video, changing signer appearance while preserving linguistic
content~\cite{xiaDiffusionModelsSign2024}, and pose-based appearance
transfer re-renders one signer's motion under another's appearance in
skeleton space~\cite{moryossefPoseBasedSignLanguage2025}. This thesis
repurposes that edit family as training-time counterfactuals, in pixel
space so that the RGB modality the rest of the pipeline consumes is
retained, and asks what the edits buy the recogniser rather than the
privacy of the signer. It differs from SDDA in pairing generative edits
with motion-locked re-renders, in asking what the intervention does to
the embedding space rather than only to task accuracy, and in testing
both across two diversity regimes.
% Framing note (2026-07-11, applied 2026-07-28): the qualification of story
% point 1 ("not well-tackled in the SL domain" rather than "unexplored")
% acknowledges Atwell/Desai/SDDA while keeping the novelty claim narrow and
% true; aligned with the SDDA scoping in the intro (pivot paragraph) and the
% Synthetic data section.

\section{Benchmarks and the diversity axis}
\label{sec:related-work-benchmarks}
As surveyed by \citet{kalinowskiMachineLearningSign2026}, a large share of
sign language recognition research has been conducted on two
studio-recorded benchmarks, Phoenix2014~\cite{kollerContinuousSignLanguage2015}
(extended to Phoenix2014T for the SLT task~\cite{camgozNeuralSignLanguage2018})
and CSL-Daily~\cite{CSL-daily}. These corpora hold recording conditions
almost constant, and on them RGB architectures were competitive
(\S\ref{sec:related-work-modality}). Their homogeneity is also their
documented weakness: \citet{jangSigningOutsideStudio2022} note the lack of
variance in setting and background, and \citet{CE-CNSL} identify the
studio setting as a bottleneck and release CE-CNSL to address it. More
recently, \citet{alyamiIsharahLargeScaleMultiScene2026b} target both the
lack of diversity and signer-independent evaluation, and their Isharah
dataset supplies the CSLR track of the SignEval 2025
Challenge~\cite{luqmanSignEval2025Challenge2025a}. As benchmarks grew more
varied, the field's preference shifted towards pose input, whose
appearance invariance is most valuable exactly where appearance varies
most.

ASL Citizen~\cite{aslcitizen} belongs to this harder lineage: it is
signer-disjoint by construction and, unlike a single-studio corpus,
aggregates recordings made by participants on their own webcams, so
lighting, framing and background vary along with the signer. It also
ships per-participant demographic metadata, which makes the composition
of its signer population inspectable rather than
assumed~\cite{atwellStudyingMitigatingBiases2024}. We adopt it as the
high-diversity setting of this thesis, and NGT744, recorded from four
signers in one studio, as the low-diversity setting: the two mark
opposite ends of the axis this thesis is about, while both evaluate
signer-independently.

\section{Input modality}
\label{sec:related-work-modality}
That shift has a clear empirical basis. On the studio benchmarks, RGB CNN
architectures such as
SlowFastSign~\cite{ahnSlowFastNetworkContinuous2023} achieved strong
performance; on the more varied ones, pose-based representations show
improved robustness at competitive
accuracy~\cite{minCloserLookSkeletonbased2025}, a trend usually
attributed to the inherent
invariance of pose representations to appearance factors such as clothing,
background, and
lighting~\cite{moryossefEvaluatingImmediateApplicability2021,
kratimenosIndependentSignLanguage2020, tranRegionAwarePoseModeling2025}.
Pose-based models, however, discard substantial
visual information present in the original RGB signal, raising the question
of whether the improved robustness of pose-based approaches stems solely
from appearance invariance, or whether RGB-based models simply fail to
capture motion cues in the presence of strong appearance signals.

The robustness attributed to pose input is also bounded by the extractor
that produces it, which is itself a learned RGB model. The standard
training and evaluation sets for pose estimation under-represent
darker-skinned individuals~\cite{lachanceCaseStudyFairness2023,
hasanHi5SyntheticData2024}, the few existing fairness audits of pose
estimators reach conclusions that flip with the choice of evaluation
set~\cite{lachanceCaseStudyFairness2023, xiangFairHumanCentricImage2025},
and on ASL Citizen the pose-based baseline shows larger skin-tone and
gender performance disparities than its RGB
counterpart~\cite{atwellStudyingMitigatingBiases2024}. There are
mechanical failure modes too: MediaPipe Holistic's hand
region-of-interest prediction breaks down on non-ideal hand
orientations~\cite{moryossefOptimizingHandRegion2024}, and both OpenPose
and MediaPipe degrade precisely where hands interact with each other or
the face, interactions that are linguistically load-bearing in
signing~\cite{moryossefEvaluatingImmediateApplicability2021}. In our own
extraction MediaPipe produced no usable feature file for $8\%$ of ASL
Citizen training videos (Table~\ref{tab:raw_feature_l2_probe}).

The appearance bias of RGB models is, however, a property of their
training data rather than of the modality itself
(\S\ref{sec:related-appearance-bias}); the same evidence shows it
survives the move to transformer backbones, which
\citet{liMitigatingEvaluatingStatic2023a} mitigate with augmentation on a
video Swin transformer. That is the question we ask of the transformer
backbones now common in SLR.

\subsection{Single-stream scope}
A natural alternative is
to combine RGB and pose in hybrid architectures, e.g.\ a two-stream
network~\cite{chenTwoStreamNetworkSign}; we do not pursue this, preferring a
single-stream RGB approach for its lower inference cost and simpler
pipeline, considerations that real-time sign language processing makes
explicit~\cite{moryossefRealTimeMultilingualSign2024}. Throughout this thesis we therefore consider \emph{single-stream}
methods only: every backbone we train, probe, or compare consumes one input
stream (RGB video or pose), and multi-stream or hybrid architectures are
out of scope on compute and inference-cost grounds.

\section{Synthetic data}
\label{sec:related-work-synthetic}
If appearance bias is inherited from training-data statistics
(\S\ref{sec:related-appearance-bias}), reshaping the training
distribution is a way to reduce it. Sign language recognition, however, lacks a counterpart
to the large, temporally rich pretraining corpora available in general video
understanding, motivating the use of \emph{synthetic} data to reshape the
training distribution directly. Knapp and Bohacek~\cite{knapp2025synthetic}
generate synthetic action video via controllable 3D Gaussian avatar pose
transfer and show gains when used as augmentation; related work renders
synthetic action data in GTA V~\cite{guoLearningVideoRepresentations2022} and
studies synthetic data for temporal alignment~\cite{SVLTA_Du_2025}. Li et
al.\ show that data generated by a text-to-video diffusion transformer can
benefit data-limited action
understanding~\cite{liRoleVideoGeneration2025}, but do not consider sign
language; Knapp and Bohacek further note that sign-language tasks are among
those that have so far been unable to exploit the full potential of
synthetic video, as generated human motion suffers from uncanny
artefacts~\cite{knapp2025synthetic}; this thesis addresses that gap.

In general vision, generative augmentation is itself an established
line: DA-Fusion edits real training images with a diffusion model and
improves few-shot classification, while cautioning that the gains are
task- and regime-dependent~\cite{trabuccoEffectiveDataAugmentation2024},
a caveat our two-regime result echoes in the sign domain.

\subsection{Sign language data augmentation approaches}
Walsh et al.\ show that sign language production can generate useful
synthetic augmentations for sign language translation, with a revealing
modality split: skeleton-level sign-stitching augmentation yields large,
consistent gains (at least 15\% across all reported metrics), whereas
photo-realistic RGB augmentation, via a GAN-based renderer (SignGAN) and
a Gaussian-splatting renderer
(SignSplat~\cite{ivashechkinSignSplatRenderingSign2025}), yields smaller
and mixed improvements, with GAN artefacts sometimes degrading
performance~\cite{walshUsingSignLanguage2025}. Diffusion-based sign
production continues this photo-realistic generation
line~\cite{SignDiff}.
% Verified against the paper (2026-07-11): pose-vs-RGB contrast confirmed,
% but the RGB methods compared are SignGAN + SignSplat — NOT diffusion.
% Story point 8 says "(Gaussian splatting, diffusion) don't work as well";
% the "diffusion" attribution is not supported by Walsh et al. — flagged
% to Piotr in the chat summary rather than silently corrected in his story.

Fu et al.'s Signer Diversity-driven Data Augmentation (SDDA) combines
diffusion-generated signer variation with a contrastive/adversarial
consistency objective and reports gains on
Phoenix2014T \cite{fuSignerDiversitydrivenData2024b}; the authors note that
their synthesized signers do not closely resemble real footage and frame
this as a limitation. We use the same consistency-objective family as one of
our backbone-fine-tuning ablations (\S\ref{sec:method-objectives}) and
return to the comparison in our Discussion (\S\ref{sec:discussion-regimes}):
SDDA's positive result is on a studio-recorded, lower-signer-diversity
corpus, closer to our low-diversity NGT744 regime than to our high-diversity ASL Citizen
regime, which we argue explains why the same idea saturates in one setting
and helps in the other.

~\citet{pengSyntheticViewAugmentation2025a} propose synthetic view augmentation using SMPL-X reconstruction
to address viewpoint generalisation, introducing new backgrounds and
viewpoints but not new signer
appearance; notably, their
synthetic viewpoints enter training only through the projected-keypoint
(pose) and rendered-depth streams of a multi-modal model, with view
augmentation of the RGB stream explicitly left to future work, whereas
this thesis targets the lack of variety in signer appearance directly, in
the RGB modality.

\subsection{Avatar and Unreal Engine approaches}
The closest precedent for our motion-locked route is Ranum's thesis on
synthetic avatar data for Dutch Sign Language
(NGT)~\cite{ranum3DAwarenessGeometry2024}. Like us, Ranum renders
synthetic avatars driven by recorded motion capture; unlike us, the
renders are evaluated in the pose modality, by extracting keypoints from
the rendered videos and training on those. In that modality the
synthetic data benefits recognition, but the thesis expects the gain to
transfer less well to the video modality, given the cartoon-like
appearance of the avatars used. Our experiments test exactly that
video-modality case: the renders enter training as RGB clips, appearance
included. 

Prior work has built gesture-producing virtual avatar systems using Unreal
Engine and motion capture~\cite{messinaModularSignLanguage2025}, and has
used Unreal-rendered virtual agents, animated from mocap and varied in
skin tone, lighting, background, and camera angle, as a scalable source
of training data for gesture recognition~\cite{loyVirtualAgentsScalable}.

What our work adds to this line is the counterfactual construction and
its controlled evaluation. Prior avatar pipelines generate synthetic
gesture or sign data at scale; none of them, to our knowledge, renders
the same captured signing under systematically varied appearance and
then asks, against matched real controls, what that variation buys. We
evaluate the renders directly as RGB training views for sign language,
compare them at matched view count against a real multiview control and
against generative edits of the same videos, and trace the effect into
the learned embedding space rather than reporting task accuracy alone.
